% =====================================================================
\chapter{Literature Review}
\label{ch:literature}

\section{Overview}
Real-time bus tracking sits at the intersection of several research
areas: vehicle telematics, mobile and participatory sensing, arrival-time
prediction, and GPS signal processing. This chapter reviews the work most
relevant to a low-cost, passenger-facing university deployment, grouped
into five strands: campus bus tracking systems
(Section~\ref{sec:lit-campus}), smartphone-based and crowdsourced
positioning (Section~\ref{sec:lit-phone}), bus arrival-time prediction
(Section~\ref{sec:lit-eta}), GPS quality and map matching
(Section~\ref{sec:lit-gps}), and open telematics platforms
(Section~\ref{sec:lit-telematics}). Section~\ref{sec:lit-gap} then
summarizes the gap that \system{} addresses.

\section{Campus and Real-Time Bus Tracking Systems}
\label{sec:lit-campus}
Student-facing bus tracking has attracted steady attention because the
need is so tangible. Sharif \emph{et al.}~[1] describe a real-time campus
bus tracking mobile application in which an on-board GPS tracker is paired
with an IoT passenger counter, so that riders see both the location of
the bus and an indication of how full it is. Premkumar \emph{et al.}~[2]
present a college bus tracking and notification system that combines GPS
tracking with a schedule view and push notifications, so that students are
actively told about the bus rather than having to watch a map
continuously.

These systems demonstrate the practical value of live tracking on a
campus fleet, and several of their ideas---push notifications, schedule
integration---also appear in \system{}. However, they share two
characteristics that limit their applicability to deployments like ours.
First, each depends on a \emph{single} on-board positioning source; if
that device fails, is powered off, or is not yet installed on a
particular vehicle, tracking for that bus simply stops. Second, their
interfaces are English-only, which is a real adoption barrier for a
student population that is more comfortable in Bangla. \system{} differs
on precisely these two axes: it ingests two independent position sources
with automatic arbitration between them, and its entire passenger
experience is bilingual.

\section{Smartphone-Based and Crowdsourced Positioning}
\label{sec:lit-phone}
A complementary research line removes dedicated hardware entirely and
uses smartphones as the position source. The seminal work of Zhou
\emph{et al.}~[3] showed that bus arrival times can be predicted from
passengers' phones alone: participatory sensing infers which passengers
are on a bus and where that bus is, without any cooperation from the
transit operator. Wepulanon \emph{et al.}~[4] develop this idea into a
real-time arrival information framework driven entirely by crowdsourced
smartphone observations, evaluated through simulation experiments.

The appeal of these approaches is obvious---zero hardware cost---but they
carry a structural weakness for a small university fleet: coverage depends
on enough contributing phones being on the bus at the right time. A
route's first trip of the morning, a lightly loaded off-peak trip, or a
trip whose riders decline to share location data all produce gaps
precisely when a waiting passenger most needs information.

\system{} deliberately \emph{inverts} the relationship between phones and
hardware. Normal operation depends on no phone at all: the primary source
is an autonomous, vehicle-fixed GPS tracker. A phone---the driver's, not
the passengers'---participates only as an optional fallback that is
selected automatically when no healthy hardware fix is available. Phone
sensing thus supplements the system instead of carrying it.

\section{Bus Arrival-Time Prediction}
\label{sec:lit-eta}
Predicting when a bus will arrive is a well-studied problem with a
spectrum of approaches. Reich \emph{et al.}~[5] survey ETA prediction
methods for public transport and organize them by the input data they
require, while Singh and Kumar~[6] review artificial-intelligence
approaches specifically for bus arrival-time prediction.

On the data-driven end of the spectrum, Ye and Thiengburanathum~[7]
demonstrate support-vector regression running on a Spark big-data
framework for city buses in Chiang Mai; Rashvand \emph{et al.}~[8],~[9]
train fully connected neural networks that predict arrival and departure
times across hundreds of routes; Jeong \emph{et al.}~[10] apply a
transformer encoder to capture long-range temporal dependence in bus
trajectories; and Google's BusTr system, described by Barnes
\emph{et al.}~[11], translates real-time traffic forecasts into bus delay
predictions at a global scale. These models achieve strong accuracy, but
they share a prerequisite that a brand-new single-campus deployment cannot
meet: large historical trajectory datasets, and in some cases external
traffic feeds, must exist before the model can be trained at all.

\system{} therefore takes the pragmatic position appropriate to a cold
start: a lightweight \emph{geometric} estimator (remaining along-route
distance divided by a rolling-average speed) that needs no training data,
paired with systematic logging of every prediction so that a data-driven
model can be trained later, once sufficient operational history has
accumulated. The estimator is described in Section~\ref{sec:design-eta}.

\section{GPS Quality, Filtering, and Map Matching}
\label{sec:lit-gps}
Raw GPS is noisy, and a public-facing tracker cannot simply display every
fix it receives. The classical treatment of reconciling noisy fixes with a
road network is the Hidden Markov Model map-matching method of Newson and
Krumm~[12], which jointly considers measurement noise and road-network
topology to recover the most probable path. HMM map matching is the right
tool when the vehicle could be anywhere on an arbitrary road network.

A university bus, however, is not anywhere: it follows a small set of
fixed, pre-defined routes. \system{} exploits this constraint to use a far
lighter pipeline---an accuracy gate, a physical-plausibility (teleport)
guard, a stationary-jitter floor, and projection of each fix onto the
known route polyline. This subset is sufficient to stabilize the displayed
position and to drive along-route ETA computation, at a small fraction of
the computational cost of full map matching
(Section~\ref{sec:design-filter}).

\section{Open Telematics Platforms}
\label{sec:lit-telematics}
Commercial GPS trackers speak a large variety of vendor-specific binary
protocols. Open-source telematics servers such as Traccar~[13] exist
precisely to absorb this diversity: Traccar implements protocol decoders
for a wide range of tracker models and exposes the decoded positions
through a REST API and WebSocket stream, along with geofencing
primitives. Building on such a platform means a project does not
reimplement low-level device communication.

\system{} adopts this approach: a self-hosted Traccar instance handles raw
device ingestion for the Concox GT06S tracker, and the entire
passenger-facing, ETA-aware, bilingual application layer is built on top
of it as an independent application server
(Section~\ref{sec:design-acquisition}).

\section{Research Gap and Positioning of This Work}
\label{sec:lit-gap}
Table~\ref{tab:lit-compare} summarizes the review. Across the reviewed
strands, three gaps recur. First, campus tracking prototypes rely on a
single position source and stop working when it fails. Second,
crowdsourced systems avoid hardware but inherit participation-dependent
coverage. Third, accurate ETA models presuppose historical data that a new
deployment does not have, and most reported systems are either
simulation-based or English-only.

\begin{table}[H]
\caption{Positioning of \system{} against the reviewed literature.}
\label{tab:lit-compare}
\centering
\small
\begin{tabular}{@{}p{4.6cm}p{4.5cm}p{4.6cm}@{}}
\toprule
\textbf{Strand} & \textbf{Typical limitation} & \textbf{\system{} response} \\
\midrule
Campus trackers [1], [2] & Single source; English-only &
Dual-source with arbitration; bilingual UX \\
Crowdsourced sensing [3], [4] & Coverage needs participants &
Autonomous fixed tracker; phone only as fallback \\
Data-driven ETA [7]--[11] & Needs large training history &
Training-free geometric ETA; logs data for future models \\
HMM map matching [12] & Heavier than needed for fixed routes &
Route-constrained filtering and projection \\
Telematics platforms [13] & Operator-facing, no passenger layer &
Passenger-centric application layer on top \\
\bottomrule
\end{tabular}
\end{table}

In contrast to single-source prototypes, simulation-only studies, and
data-hungry prediction models, \system{} contributes a \emph{deployed},
low-cost system that combines source resilience, a training-data-free ETA
estimator, and a localized, edge-condition-aware passenger experience,
evaluated on a real campus fleet over seven weeks. To the best of our
knowledge, no previously reported university bus tracking system combines
these properties.

\section{Summary}
This chapter surveyed five strands of related work and identified the
recurring gaps: single-source fragility, participation-dependent
coverage, training-data prerequisites, and the absence of localized,
passenger-centric deployments. The next chapter translates the resulting
opportunity into concrete requirements for \system{}.
